\documentclass[12pt,a4paper]{article}
\usepackage{amsmath}
\usepackage{amssymb}
\usepackage{bm}
\usepackage[utf-8]{inputenc}
\usepackage{xeCJK}

\setCJKmainfont{Hiragino Sans GB}

\title{京都大学大学院 数理工学専攻}
\author{凸最適化 過去問}
\date{H28 (OR)}

\begin{document}

\maketitle

関数 $f:\mathbb{R}^n \to \mathbb{R}$ は 2 回連続的に微分可能な関数とし，$a$ は $0$ でない $n$ 次元ベクトルとする。
次の非線形計画問題を考える。

\begin{align*}
\text{P:} \quad & \text{Minimize} \quad f(x) \\
& \text{subject to} \quad a^T x = 0
\end{align*}

ただし，$^T$ はベクトルの転置を表す。$x^*$ を問題 P の大域的最適解とする。
さらに，次の非線形計画問題を考える。

\begin{align*}
\text{P}(k): \quad & \text{Minimize} \quad f_k(x) \\
& \text{subject to} \quad (x - x^*)^T(x - x^*) \leq 1
\end{align*}

ただし，$k$ は非負の整数であり，$f_k:\mathbb{R}^n \to \mathbb{R}$ は以下に定義された関数である。

\[f_k(x) = f(x) + \frac{k}{2}(a^T x)^2 + \frac{1}{2}(x - x^*)^T(x - x^*)\]

問題 P(k) の大域的最適解を $x^k$ とする。さらに，
\[\lim_{k \to \infty} x^k = \bar{x}, \qquad \lim_{k \to \infty} k(a^T x^k) = \bar{\lambda}
\]
と仮定する。以下の問いに答えよ。

\section*{(i)}

任意の非負の整数 $k$ に対して，$f_k(x^k) \leq f(x^*)$ が成り立つことを示せ。

\section*{(ii)}

$a^T \bar{x} = 0, \ \bar{x} = x^*$ となることを示せ。

\section*{(iii)}

問題 P(k) のカルーシュ・キューン・タッカー（Karush-Kuhn-Tucker）条件を書け。

\section*{(iv)}

十分大きな $k$ に対して，$\nabla f_k(x^k) = 0$ となることを示せ。

\section*{(v)}

$\nabla f(x^*) + \bar{\lambda} a = 0$ となることを示せ。

\end{document}
